\section{Device Results}
\label{sec:device}

\subsection{Calibration quality}

\begin{figure*}[p]
\centering\includegraphics[height=0.9\textheight]{panel02.png}
\caption{Calibration against published measurements of a comparable gate-all-around
FeFET. (a) Overlay in the operating region, with the memory window and both threshold
voltages computed from the plotted curves at the stated constant-current criterion. A
single free width scalar, fitted on the erased branch's saturation, is applied to both
simulated branches; the vertical axis is labelled accordingly. (b) The same data with the
window and the current level divided out --- each curve against its own gate overdrive,
normalized to its own on-current --- so that only the turn-on shape remains. (c) Residual
in decades against the measurement, with a $\pm\num{0.3}$ decade band. Computed values:
\fignote{B1}; \fignote{B3}; \fignote{B5}.}
\label{fig:panel2}
\end{figure*}

Figure~\ref{fig:panel2} shows the calibrated deck against the measurement. On the overlay
(Fig.~\ref{fig:panel2}(a)) the simulated branches follow the measured points through turn-on and into
saturation on both polarization states, and the memory window extracted from the plotted
curves at a constant current of \SI{1e-8}{\ampere} per nanosheet is \SI{1.296}{\volt},
with threshold voltages of \SI{+0.362}{\volt} erased and \SI{-0.934}{\volt} programmed.
Figure~\ref{fig:panel2}(b) removes both the window and the absolute current level by plotting each curve
against its own gate overdrive and normalized to its own on-current: the turn-on shape
still tracks the measurement, which is the part of the agreement that the fitted width
scalar cannot manufacture. Figure~\ref{fig:panel2}(c) gives the residual, which stays inside
\num{0.3} decades over most of the operating region; the RMS residual is
\num{0.30} decades on the programmed branch and \num{0.48} on the erased branch,
the larger erased value coming from the deep subthreshold points where the measurement
itself is near its noise floor.

\subsection{The ferroelectric, and what the gate stack actually sees}

\begin{figure*}[t]
\centering\includegraphics[width=\textwidth]{panel03.png}
\caption{The ferroelectric material against what the gate stack reaches. (a) Standalone
metal--ferroelectric--metal capacitor at three sweep amplitudes. Twice the coercive
voltage does not saturate the loop; the saturated amplitude reproduces the calibrated
$P_s$, $|P_r|$ and $F_c$, and the unsaturated amplitudes are retained as the control.
(b) The loop measured through the gate stack, sign-corrected, against the same material's
coercive field. A \SI{\pm4}{\volt} gate sweep reaches about half of $F_c$: this is a deep
minor loop of the stack, not the material loop. Computed values: \fignote{C1};
\fignote{C2}.}
\label{fig:panel3}
\end{figure*}

Figure~\ref{fig:panel3} contains the paper's clearest single correction. Figure~\ref{fig:panel3}(a) is a standalone
metal--ferroelectric--metal capacitor swept at three amplitudes. At
\SI{\pm1.96}{\volt}, which is twice the coercive voltage
$2 F_c T_\mathrm{fe}$ and the amplitude a reader would reasonably choose, the loop is
\emph{not} saturated. Saturation needs roughly five times the coercive field, and there
the extracted parameters are $P_s=\SI{39.99}{\micro\coulomb\per\square\centi\metre}$,
$|P_r|=\SI{31.99}{\micro\coulomb\per\square\centi\metre}$ and
$F_c=\SI{1.400}{\mega\volt\per\centi\metre}$, against the calibration targets of
\num{40}, \num{32} and \num{1.4}. The unsaturated amplitudes are kept in the figure as
the control: they are what makes the saturated loop evidence rather than an assertion.

Figure~\ref{fig:panel3}(b) is the loop measured through the gate stack over the same material. It is a
deep minor loop. A \SI{\pm4}{\volt} gate sweep delivers about
\SI{0.77}{\mega\volt\per\centi\metre} to the HZO, roughly \SI{55}{\percent} of $F_c$, so
what the gate can switch is set by the metal--ferroelectric--insulator--semiconductor
voltage divider rather than by the material. Reading Fig.~\ref{fig:panel3}(b) as an HZO material loop
(which an earlier version of this work did) understates $P_r$ by almost an order of
magnitude and misattributes a stack property to the film.

\subsection{Retained characteristics}

\begin{figure*}[t]
\centering\includegraphics[width=\textwidth]{panel04.png}
\caption{DC characteristics. (a) Retained transfer characteristic, both states from the
same run and read after settling. Threshold markers are at the stated constant-current
criterion, the subthreshold slopes are fitted over at least two decades starting a decade
above the branch minimum, and the ambipolar minima are labelled rather than cropped.
(b) Output characteristics for the erased state, the programmed state and intermediate
ladder stops; the inset is the low-drain-bias region, with the computed linear-region
correlation given below the axes. Computed values: \fignote{D1}; \fignote{D2}.}
\label{fig:panel4}
\end{figure*}

Figure~\ref{fig:panel4}(a) is the retained transfer characteristic. After a \SI{\pm2}{\volt} write and
settling, the two states are separated by \SI{0.387}{\volt} at a constant current of
\SI{1e-2}{\micro\ampere\per\micro\metre} and by $5410\times$ in drain current at
$V_G=0$. The subthreshold slope, fitted over more than two decades starting a decade
above the branch minimum, is \SI{63.5}{\milli\volt\per\text{dec}} on the programmed
branch. The ambipolar minima at $V_G\approx\SI{-0.25}{\volt}$ (erased) and
$\SI{-0.55}{\volt}$ (programmed) are labelled rather than cropped: they are real, they
set where the device can be read, and Fig.~\ref{fig:panel11}(b) turns them into a design constraint.

Figure~\ref{fig:panel4}(b) shows that every analog state behaves as an ordinary transistor. The output
characteristic is ohmic at low drain bias for the erased state, the programmed state and
the intermediate ladder stops alike, with linear-region correlation coefficients of
\numrange{0.958}{0.998}, and saturates cleanly. Stored polarization changes the channel
conductance without changing the character of the device.

\paragraph*{Drain-induced barrier lowering}
Extracted at a fixed constant current of
\SI{6.338e-3}{\micro\ampere\per\micro\metre}, the erased branch gives
\SI{59.5}{\milli\volt\per\volt}: an unremarkable short-channel value for a
\SI{100}{\nano\metre} gate on a \SI{5}{\nano\metre} body. The programmed branch gives
\SI{309}{\milli\volt\per\volt}, five times larger and not explicable as electrostatic
barrier lowering at this geometry. The plausible mechanism is the drain field acting on
the ferroelectric itself, a drain-induced polarization change rather than a short-channel
effect. \emph{This run does not establish that.} Settling it needs the polarization
probed at the drain end against $V_{DS}$, or a non-ferroelectric control, and neither was
run. The erased value is reported as the device's DIBL; the programmed value is reported
as an open question rather than omitted.

\subsection{Write transients and energy}

\begin{figure*}[p]
\centering\includegraphics[height=0.9\textheight]{panel05.png}
\caption{Synaptic behaviour. (a) The potentiation ladder with its measured cycle-to-cycle
band. Filled markers are separable at three standard deviations, open markers are not; the
levels and the noise come from the same eleven-repeat run. (b) The conventional
exponential-saturation nonlinearity fit, with the residual shown as an inset: it is
systematic and one-sided because the train is sigmoidal, so the \numrange{10}{90} percent
pulse span is quoted alongside the fitted parameter. (c) Three potentiation--depression
cycles on a log axis. The window lost between cycles is lost at the floor, not the
ceiling. This is write repeatability, not endurance. Computed values: \fignote{F1};
\fignote{F2}; \fignote{F5}.}
\label{fig:panel5}
\end{figure*}

\begin{figure*}[p]
\centering\includegraphics[height=0.9\textheight]{panel06.png}
\caption{Write transients and energy, all from one fifteen-pulse train. (a) Polarization
and gate charge stitched onto one time axis, with the write windows shaded.
(b) Per-pulse polarization transients, colour-graded by pulse index, and the amount
switched per pulse: switching peaks in the middle of the train and falls as the film
saturates. (c) Switching energy per pulse and its cumulative total, taken from the
solver's gate-charge trace rather than from a $V_G I_G$ integral, which at this sampling
interval would measure the sampler. Computed values: \fignote{E1}; \fignote{E2};
\fignote{E6}.}
\label{fig:panel6}
\end{figure*}

Figure~\ref{fig:panel6} follows one train of fifteen \SI{2}{\volt}, \SI{100}{\nano\second} pulses.
Figure~\ref{fig:panel6}(a) stitches the write and read segments onto one time axis: the polarization
walks down a staircase, one step per pulse, and between pulses the gate conduction
current sits at a median \SI{26.7}{\pico\ampere}, so the \SI{1}{\nano\metre} interfacial oxide
does not bleed the stored state away between writes.

Figure~\ref{fig:panel6}(b) overlays the per-pulse polarization transients. The amount switched per pulse
rises through the first third of the train, peaks, and then falls as the ferroelectric
saturates. That kinetic saturation (not any limit of the read) is why the
conductance ladder compresses at the top, and it is the mechanism behind the level
crowding reported in Section~\ref{sec:levels}.

Figure~\ref{fig:panel6}(c) gives the measured energy. Switching costs \SI{8.2}{\atto\joule} per pulse on
average, peaking at \SI{11.0}{\atto\joule} in the middle of the train, and the complete
fifteen-level write costs \SI{0.123}{\femto\joule}.

\subsection{Synaptic behaviour}

Figure~\ref{fig:panel5}(a) is the potentiation ladder with its measured cycle-to-cycle band, discussed in
Section~\ref{sec:levels}. Figure~\ref{fig:panel5}(b) fits the conventional exponential-saturation
nonlinearity model to the normalized ladder. It fits poorly, and the figure says so: the
residual is systematic and one-sided because this device's response is sigmoidal (an
incubation of two to three pulses, then a fast swing) rather than saturating from the
first pulse. The nonlinearity parameter is reported alongside the \numrange{10}{90}
percent pulse span, which describes this shape and the single parameter does not.

Figure~\ref{fig:panel5}(c) closes the loop. Three alternating potentiation and depression trains give
per-cycle windows of $11650\times$, $4185\times$ and $3577\times$, and the depression
endpoint returns to \SI{96.2}{\percent} of the erased baseline. The window that is lost
between cycles is lost almost entirely at the \emph{floor}: the erased minimum rises from
\num{0.0040} to \SI{0.0121}{\micro\ampere\per\micro\metre} across the three cycles while
the potentiated ceiling falls only from \num{46.9} to
\SI{43.4}{\micro\ampere\per\micro\metre}. A single erase pulse does not fully switch a
programmed device, because erase is opposed by the depolarization field.

\emph{This is write repeatability, not endurance.} The Preisach model has no fatigue,
wake-up or imprint term, so it is structurally incapable of producing an endurance
result: cycling it further would return a flat line because the equations say so. No
endurance figure is presented and none should be inferred.

Temperature is reported in the supplementary set and is not monotonic: at the top of the
ladder \SI{350}{\kelvin} sits about six times \emph{below} \SI{300}{\kelvin} while
\SI{250}{\kelvin} sits slightly above it. Three temperatures, non-monotonic, do not
support an Arrhenius fit, and none is presented.

\subsection{Analog depth: eight levels, or sixteen}
\label{sec:levels}

Figure~\ref{fig:panel5}(a) is the central device measurement of this paper. Eleven complete repeats of
the same fifteen-pulse train on one device --- a genuine cycle-to-cycle distribution
rather than a corner envelope --- give a median standard deviation of \num{0.0215}
decades, with the per-level value falling from \num{0.239} decades at the bottom of the
ladder to \num{0.0131} at the top. The run was stopped by a two-hour job limit after
eleven of the twenty queued repeats, so the standard deviation itself is known to about
\SI{22}{\percent}, and that repeat count is quoted wherever the number is.

At a three-sigma separation criterion, \textbf{eight of the fifteen stops are separable}.
The failures are all at the top: levels nine to fifteen sit inside \num{0.114} decades
because the potentiation curve saturates, while the noise stays roughly flat, so at level
fifteen the spacing is \num{0.0085} decades against a standard deviation of \num{0.0131}
--- the levels are closer together than the noise.

The measured range is \num{3.79} decades wide, and the train's own stops are not the only
places a level can be put. Placing targets greedily from the bottom, each three standard
deviations above the last, the same range and the same measured noise support about
\textbf{sixteen levels}. Reaching them requires write-verify circuitry and per-write
iteration.

\textbf{Both numbers are reported, with their protocols attached, and the network results
in Section~\ref{sec:network} assume a write-verify array.} That is not a concession made
to rescue a level count: Section~\ref{sec:variability} shows that write-verify is
unavoidable on this device regardless of how many levels are claimed, because without
per-device calibration the deployed accuracy collapses to \num{0.243}. The system pays
for write-verify either way, so the higher level count costs nothing extra.

\subsection{Retention and read disturb}

Figure~\ref{fig:panel9} is built around one rule: measure after the state has settled. Figure~\ref{fig:panel9}(a) shows
why. A hold immediately after a write is dominated by a depolarization transient --- a
roughly ninefold overshoot decaying away over about five time constants --- and only the
plateau afterwards is retention. Fitting a single power law across both phases returns
total data loss at ten years from a signal that is provably flat, and an unphysical
negative activation energy. That is not a hypothetical: it is what this work reported
before the two phases were separated.

Figure~\ref{fig:panel9}(b) gives the retention itself, out to \SI{10}{\milli\second}. After settling at
about \SI{50}{\micro\second}, the programmed state at \SI{300}{\kelvin} is flat to
\SI{0.62}{\percent} across \num{2.3} decades; at \SI{400}{\kelvin} it is flat to
\SI{0.96}{\percent} across \num{2.4} decades; and a mid-ladder analog state, LTP level
eight, is flat to \SI{0.62}{\percent} --- the same bound as the fully programmed state,
which is the case reviewers most often doubt.

No ten-year extrapolation is given. The measurement spans \SI{10}{\milli\second}; ten
years is eleven decades beyond it, and the plateau is flat to the solver's own
resolution, so an extrapolated number would be reporting a fit's noise floor as physics.
What the data supports is a bound: no decay resolvable above about \SI{1}{\percent} over
the measured decades. No activation energy is extracted either --- both temperatures
reach flat plateaus, the plateau \emph{levels} differ mainly through the temperature
dependence of the drain current rather than through polarization loss, and the model
contains no thermally activated retention term.

Figure~\ref{fig:panel9}(c) is read disturb. Measured from the settled state, the retained current changes
by \SI{0.0000}{\percent} over \num{9.9e5} equivalent reads at $V_\mathrm{read}=0$. This
is a strong result and it has a simple mechanism: at zero gate bias the read does no work
on the ferroelectric, so there is nothing to disturb it. The same run measured from the
immediate post-write value would read \SI{-89.1}{\percent}, which is the depolarization
transient and not disturb; both are drawn so the distinction is visible rather than
asserted.

\subsection{Integrate-and-fire operation}

\begin{figure*}[p]
\centering\includegraphics[height=0.9\textheight]{panel07.png}
\caption{Integrate-and-fire operation. (a) Read current against pulse number at
$V_G=0$, crossing a fixed threshold on the ninth pulse at $18.4\times$ the rest current,
measured against this run's own baseline. (b) Membrane-state accumulation for five drive
amplitudes, measured on the polarization probe during the write; pulse amplitude sets the
integration rate. (c) Gate drive, polarization and read current on one time axis. With the
depolarization time constant used here the device latches after firing, so this is a
single-shot event rather than a repeating leaky cycle. Computed values: \fignote{G1};
\fignote{G2}; \fignote{G3}.}
\label{fig:panel7}
\end{figure*}

Figure~\ref{fig:panel7} shows the device used as an integrator. In Fig.~\ref{fig:panel7}(a), successive identical
pulses raise the zero-bias read current monotonically until it crosses a fixed threshold
of \SI{0.1}{\micro\ampere\per\micro\metre} on the ninth pulse, at a contrast of
$18.4\times$ the rest current --- taken against that run's own pre-train baseline.
Against a different run's baseline the same measurement reads $30.6\times$, which is the
error the same-run rule in Section~\ref{sec:rules} exists to prevent.

Figure~\ref{fig:panel7}(b) measures the state variable directly. Across five programming amplitudes from
\SIrange{1.5}{3.5}{\volt}, the polarization accumulated over a nine-pulse train reaches a
fixed threshold in seven, six, five, four and four pulses respectively: pulse amplitude
sets the integration rate, which is the transfer function an integrate-and-fire element
needs. Figure~\ref{fig:panel7}(c) puts the three traces on one time axis --- gate drive, polarization,
read current --- and makes the point that the polarization \emph{is} the state variable.

Two limitations belong with this panel. First, with the depolarization time constant used
here the device holds its state after firing rather than leaking back, so what is
demonstrated is a single-shot integrate-and-fire event and not a repeating leaky cycle.
Second, this project ran two dedicated count-to-fire sweeps, 224 files in total, and
neither is usable. Both write with \SI{20}{\nano\second} pulses, and the programming
study in Fig.~\ref{fig:panel8} already establishes that pulses of \SI{30}{\nano\second} and shorter do
not program at any amplitude; the polarization traces agree, moving by about
\SI{0.01}{\percent} of $P_r$ across twelve pulses. Several amplitude families also read
at $V_G=\SI{-0.5}{\volt}$, inside the gate-induced-drain-leakage branch, where the
current is set by band-to-band generation at the drain edge and is nearly insensitive to
the stored polarization. Those runs completed without error and produced flat, plausible
looking ladders. They are reported here rather than quietly dropped, because a run that
completes is not a result.

\subsection{Design space}

\begin{figure*}[p]
\centering\includegraphics[height=0.9\textheight]{panel08.png}
\caption{Design space. (a) The coordinate-descent trajectory, six turns, with the value
chosen at each turn. (b) Interfacial oxide thickness: retained window against off-current,
in both evaluation contexts. Computed values: \fignote{I1}; \fignote{I2}.}
\label{fig:panel8}
\end{figure*}

\begin{figure*}[p]
\centering\includegraphics[height=0.9\textheight]{panel11.png}
\caption{(a) Ferroelectric thickness as a Pareto choice --- window
against operating voltage, marker area proportional to gate charge. (b) Read bias: the
true retained window and the rest current against $V_\mathrm{read}$, with the leakage
region shaded, and the superseded virgin-referenced dynamic range overlaid as a dashed
line to document the disagreement. Computed values: \fignote{I3}.}
\label{fig:panel11}
\end{figure*}

\begin{figure*}[p]
\centering\includegraphics[height=0.9\textheight]{panel09.png}
\caption{Retention and read disturb, all measured after settling. (a) A short hold, with
the depolarization transient and the plateau separated and the wait time stated.
(b) Long retention at two temperatures, for the fully programmed state and for a
mid-ladder analog state, normalized to each plateau. No ten-year extrapolation is drawn
and no activation energy is extracted; a bound over the measured decades is reported
instead. (c) Read disturb against equivalent read count at $V_\mathrm{read}=0$. The
transient is shown in grey so that the distinction between it and disturb is visible
rather than asserted. Computed values: \fignote{H1}; \fignote{H4}; \fignote{H8}.}
\label{fig:panel9}
\end{figure*}

Figure~\ref{fig:panel8} collects the parameter study. Figure~\ref{fig:panel8}(a) is the coordinate-descent trajectory:
six turns, four orders of magnitude of improvement in the retained window, and the last
two turns changing nothing, which is what termination looks like.

Figure~\ref{fig:panel8}(b) is the interfacial oxide. Thinning it buys window and costs off-current, and
the trade has the same shape whether each thickness is evaluated at a fixed
\SI{4}{\volt} write or at its own sub-coercive operating point. One nanometre yields
about seventy-five times the window of \SI{1.5}{\nano\metre}; below that the off-current
penalty takes over.

Figure~\ref{fig:panel11}(a) is the ferroelectric thickness, and it is a Pareto choice rather than an
optimum. Seven nanometres gives the largest retained window of the set, at the lowest
operating voltage in the set and the smallest gate charge --- so on this device the
window-optimal thickness also happens to be the cheapest to write, which is not
guaranteed and is the reason the figure plots all three quantities together.

Figure~\ref{fig:panel11}(b) is the read bias. The retained window is largest near $V_G=0$ while the rest
current is at its floor; at negative read bias the gate-induced drain leakage reopens and
the window collapses. The superseded virgin-referenced dynamic range from an earlier
version of this work is overlaid as a dashed line: it peaks in the leakage region,
because that is where the leakage current is largest, not where the stored state is most
readable. Plotting the two together documents the correction rather than hiding it, and
it is the same conclusion the unusable count-to-fire runs reach from the other direction.

\subsection{Comparison with a planar ablation}

\begin{figure*}[t]
\centering\includegraphics[width=\textwidth]{panel10.png}
\caption{Comparison and benchmark. (a) The nanosheet against a planar ablation of the
identical stack, on the axes a synapse is judged by; each axis is normalized to the better
of the two devices and the level axis uses one shared separation criterion. The planar
device wins the raw window and loses the application-relevant axes. (b) Analog level count
across a survey of fifteen recent published ferroelectric analog memories, with this work
marked at both its open-loop and its write-verify count. Every plotted count was confirmed
against its own paper's text; none of those papers states a separation criterion behind
it. Computed values: \fignote{J1}; \fignote{J5}.}
\label{fig:panel10}
\end{figure*}

Figure~\ref{fig:panel10}(a) compares the nanosheet against a planar device built from the identical
stack with the bottom gate removed --- an ablation, so gate control is the only variable.
The planar device wins the headline number outright: its raw memory window is
\SI{1.42}{\volt} against \SI{0.387}{\volt}, and its retained on/off ratio is
$1.42\times10^4$ against $5.41\times10^3$. It loses the properties this application
needs. Under one shared separation criterion the nanosheet resolves fourteen ladder stops
against the planar device's ten, it programs at \SI{2.0}{\volt} against
\SI{2.3}{\volt}, and it holds a steeper subthreshold slope. Retention is deliberately not
an axis of the comparison: the planar figure is a \SI{100}{\micro\second} hold, and on
the nanosheet \SI{100}{\micro\second} lies inside its own settling transient, so the two
numbers do not measure the same thing.

\subsection{Reported level counts are not measured against a common criterion}

Figure~\ref{fig:panel10}(b) places this device on a survey of fifteen recent ferroelectric
analog-memory papers. The comparison we wanted to draw --- analog levels against energy per
programming pulse --- cannot be drawn: only two of the fifteen state an energy per
programming pulse at all. Plotting it anyway would have put a single comparable marker on
the page and let an empty axis imply a lead that was never measured, so the figure is drawn
on the quantity the field does report, and the energy comparison is made here, in text,
where two points can be quoted as two points. This device switches at
\SI{8.2}{\atto\joule} per pulse at \SI{2.0}{\volt}; the two published values we could
verify are \SI{48}{\atto\joule} and \SI{0.8}{\femto\joule}.

On level count this device is not the leader and the figure says so. Ten of the fifteen
papers state a count, and those run from sixteen to a hundred and forty against this
work's eight open-loop and about sixteen under write-verify. Two things qualify that gap.
The counts are not measured under a common criterion --- each is whatever its own paper
called distinguishable, and this work's eight is a conservative open-loop count under its
own measured cycle-to-cycle noise rather than a best case. More importantly, the network
results in Section~\ref{sec:network} show that level count does not predict accuracy:
placement does. A device reporting a hundred and forty levels has not thereby shown
that a classifier deployed on it will work.

Every point is traceable through \texttt{literature\_benchmark.csv}, which carries the
digital object identifier and the sentence each number was read from. Six rows were
originally recorded without a confirmed count; five have since been resolved against the
source text and one paper states no count at all, so it carries no point on the figure.
One of those five moved against this work: the highest count in the survey was recorded as
a hundred and twenty-eight, inferred from a claim of more than seven bits, and the paper's
own text gives a measured hundred and forty. Not one of the five states a separability
criterion or a cycle-to-cycle statistic behind its count. The audit is in
\texttt{literature\_benchmark\_notes.md}.
